\chapter{Strands, complements and dual-state hypotheses}

Read the DNA strand $5'$--ACGA--$3'$. Its paired strand, read in its own
$5'$-to-$3'$ direction, is TCGT, not TGCT and not AGCA. Suppose a model emits a
score at every base. When the same duplex is presented from the other strand,
should the scores stay fixed, reverse position, exchange labels, or change sign?
The answer depends on the task. A biological relationship does not determine
the output representation for us.

Chapter 18 introduced strand-aware hypotheses; Chapter 20 separated memory
representations and clocks. We now construct a precise dual-direction state
operator, prove its transformation law and test its implementation. We also
prove a limitation: demanding exact reversal symmetry from an ordinary
left-causal sequence-to-sequence map can eliminate nonlocal context. DOGMA
remains a proposed non-Transformer research direction. The reference below
tests semantics, not trained genomic capability.

\section{Choose coordinates before choosing an architecture}

Use channel order $(A,C,G,T)$ and a record matrix $X\in\mathbb R^{T\times4}$.
For canonical DNA, each row is one-hot. The complement permutation is
$c=(3,2,1,0)$ in zero-based indexing. Let $P$ be its $4\times4$ permutation
matrix and $R_T$ reverse the $T$ positions. Then
\begin{equation}
 \mathcal R X=R_TXP,\qquad P^2=I,\quad R_T^2=I,\quad
 \mathcal R^2X=X .
 \label{eq:rc}
\end{equation}
Position reversal and channel complement commute because they act on different
axes. Complement alone describes the aligned partner in the original
left-to-right drawing; reversal is necessary to read that partner from its
own $5'$ end.

\Plate{01-duplex}{The physical duplex fixes the coordinate transformation.
Continuous backbone arrows point toward $3'$; short cross-strand marks denote
base pairing. Reading the lower strand right to left gives TCGT. The tensor
operation reverses rows and permutes base columns; it does not chemically
reverse a molecule.}{fig:duplex}

\Listing{rc}

The helper accepts finite real matrices with four channels so that continuous
derivatives can be tested. The canonical encoder accepts only uppercase
ACGT strings. It rejects N and lowercase input rather than silently assigning
an ambiguity policy. A production tokenizer must separately define ambiguous
bases, masking symbols, padding and record separators. Each extra symbol needs
a specified transformation; a row of zeros is not automatically a nucleotide.

For zero-based coordinates, a base at $i$ maps to $T-1-i$; a half-open interval
$[a,b)$ maps to $[T-b,T-a)$. Applying the interval rule twice returns the
original interval. A reference allele and its alternate allele must also be
complemented when transforming a variant. Reversing an array of scores without
transforming its coordinate metadata is a biological annotation error even if
the tensor dimensions agree.

\section{Invariance and equivariance answer different questions}

A record-level scalar may be invariant:
$f(\mathcal RX)=f(X)$. A per-position unstranded score should instead obey
$f(\mathcal RX)=R_Tf(X)$. A pair of strand-specific tracks may reverse position
\emph{and} swap its two labels. A four-base probability vector requires a
complement permutation as well. None of these actions is interchangeable.

In general, choose an output action $\rho$ with $\rho^2=I$ and require
\begin{equation}
 f(\mathcal RX)=\rho f(X).
 \label{eq:equiv}
\end{equation}
This is \Term{equivariance}. Invariance is the special case $\rho=I$.
It is a structural constraint, not a claim that changing orientation cannot
matter biologically. An oriented transcription task carries strand labels.
An assay, genome coordinate frame or supplied directional context can break
the symmetry unless that context is transformed too.

There is a generic full-record construction. Let the outputs lie in a vector
space and let $\rho$ be a linear involution, as all the permutation actions
above are. For any deterministic map $g$ into that space, define
\begin{equation}
 \mathcal Sg(X)=\tfrac12\{g(X)+\rho g(\mathcal RX)\}.
 \label{eq:symmetrize}
\end{equation}
Then $\mathcal Sg(\mathcal RX)=\rho\mathcal Sg(X)$ by substitution and
$\rho^2=I$. Applying $\mathcal S$ again changes nothing: it is a projection
onto equivariant maps. This construction is useful, but it evaluates $g$ on
both orientations and can discard orientation-sensitive distinctions.
For classification, averaging probabilities and averaging logits are also
different operations, even when each can be made equivariant.

Data augmentation is not this algebraic guarantee. It presents both
orientations during training, encouraging rather than enforcing a relationship.
Finite data, optimization, regularization and stochastic inference can still
produce discrepancies. A symmetry test should measure the actual deployed
operator under its actual numerical and randomness settings.

\section{Lift a causal operator into paired directional states}

Let $F_\theta:\mathbb R^{T\times4}\to\mathbb R^{T\times D}$ be a deterministic
sequence operator, shared between two evaluations. It need not itself be
complement equivariant. Define the aligned dual representation
\begin{equation}
 E_\theta(X)=
 \left[F_\theta(X)\ \middle|\ R_TF_\theta(\mathcal RX)\right]
 \in\mathbb R^{T\times2D}.
 \label{eq:lift}
\end{equation}
The second result must be reversed \emph{back} into the original coordinates.
Its hidden channels are arbitrary learned features, not nucleotide identities;
we do not apply $P$ to them. Instead define an explicit paired hidden action
\begin{equation}
 \mathcal J[U\mid V]=[R_TV\mid R_TU].
 \label{eq:pairedaction}
\end{equation}
It reverses positions and swaps feature blocks, with $\mathcal J^2=I$.

\begin{proposition}[Shared dual lifting]
For every input in the operator's domain,
$E_\theta(\mathcal RX)=\mathcal J E_\theta(X)$.
\end{proposition}
\begin{proof}
Using Equation~\ref{eq:rc}, the left side is
$[F_\theta(\mathcal RX)\mid R_TF_\theta(X)]$.
Applying $\mathcal J$ to Equation~\ref{eq:lift} gives exactly these blocks,
because $R_T^2=I$. No special recurrence equation or learned weight value
was needed. Shared parameters and compatible boundary conditions were needed.
\end{proof}

\Plate{02-lift}{The two evaluations share one parameter set but traverse
opposite orientations. Reverse-pass features are realigned before concatenation.
Changing the input orientation exchanges the roles of the two halves. Hidden
feature identity is preserved within each half; base complement and feature-block
swap are separate operations.}{fig:lift}

Choose a deliberately transparent causal reference:
\begin{equation}
 h_t=a\odot h_{t-1}+x_tW,\qquad h_0=0,
 \quad W\in\mathbb R^{4\times D},\quad a\in[0,1)^D.
 \label{eq:recurrence}
\end{equation}
Here $\odot$ denotes componentwise multiplication; rows are feature vectors.
Unlike Chapter 20's convex-average bank, this recurrence has no $(1-a)$
write factor. It sums decayed projected inputs. For bounded
$|(x_tW)_d|\leq M_d$, its zero-state magnitude is at most
$M_d(1-a_d^t)/(1-a_d)$, not $M_d$. The change in normalization is explicit.

\Needspace{29\baselineskip}
\Listing{causal}

The code validates shape, dtype, device, finiteness and retention range.
It never updates caller-owned carry in place. Empty chunks produce no output
rows and preserve carry. These checks favor inspectability, not accelerator
throughput. The returned history uses $O(TD)$ storage; a streaming forward
pass can retain only $O(D)$ state if intermediate outputs are consumed.

\Listing{dual}

For $D=1$, $W=(1,2,3,4)^\top$, $a=1/2$ and ACGA, the forward projection
sequence is $(1,2,3,1)$. The opposite strand TCGT projects to $(4,2,3,4)$.
Applying the recurrence gives $(4,4,5,6.5)$ on that strand; realignment yields
$(6.5,5,4,4)$. The generated trace checks the entire coordinate convention:
\begin{center}\input{\Generated/trace.tex}\end{center}
Adding the halves produces a position-equivariant unstranded representation.
It does not mean that the two halves are redundant: their difference carries
information that changes sign when the branch labels swap.

\section{A readout must respect the hidden action}

Write a paired feature at one position as $(u,v)\in\mathbb R^D\times\mathbb R^D$.
A general linear paired readout has blocks
$M=\left(\begin{smallmatrix}A&B\\C&D_0\end{smallmatrix}\right)$.
Let $J$ swap the two input blocks and $J'$ swap the two output blocks.
Commuting with the actions requires
\begin{equation}
 MJ=J'M .
\end{equation}
Multiplying both sides blockwise gives $A=D_0$ and $B=C$. Thus the tied form is
\begin{equation}
 M=\begin{pmatrix}U&V\\V&U\end{pmatrix},
 \qquad b=\begin{pmatrix}b_0\\b_0\end{pmatrix}
 \label{eq:tied}
\end{equation}
if an affine bias is included. The proof does not say that an arbitrary
linear layer is safe just because its input is equivariant.

\Plate{03-head}{A tied paired readout commutes with branch exchange.
Diagonal blocks share $U$ and off-diagonal blocks share $V$. The sum channel is
unchanged under exchange; the difference channel changes sign. Keeping both
is useful when output labels are strand-specific.}{fig:head}

\Listing{strand_head}

Our reference implements two scalar tracks without bias, with
$y^+=u^\top U+v^\top V$ and $y^-=u^\top V+v^\top U$.
Across a whole sequence the outputs reverse position and swap tracks.
An unstranded per-position head can use $u+v$; averaging that head over a
nonempty record makes a record-invariant summary. Empty-record pooling needs
an explicit policy and is not implemented here.

This suggests a systematic architectural audit. A pointwise nonlinearity
shared across coordinates commutes with a permutation. An arbitrary
normalization with untied affine parameters may not. Absolute positional
embeddings, branch-specific dropout masks, independent branch parameters and
orientation-dependent initial states can all break the proof. At inference,
disabled dropout is deterministic; at training, a claim of pathwise symmetry
needs correspondingly transformed randomness, not merely equal mask
distributions.

\section{Why the reverse pass cannot be ordinary causal decoding}

If $F_\theta$ is causal, its feature at position $t$ depends on $x_1,\ldots,x_t$.
The aligned second half in Equation~\ref{eq:lift} depends on
$x_t,\ldots,x_T$. Together they see the entire record. Reversing a known
record is therefore an offline operation, not a way to obtain future context
from a prefix.

\Plate{04-causality}{At a selected position, the forward branch reads the
prefix and the aligned reverse branch reads the suffix. Their overlap is the
current token; their union is the full record. The second branch is valid for
offline inference but would leak future information into a left-to-right
next-token predictor.}{fig:causality}

\begin{proposition}[A causality--reversal limitation]
Fix a record length $T$ and a full product input domain. Suppose
$f(X)_t$ is left-causal, and $f$ is equivariant under reversal with invertible
pointwise input/output channel permutations. Then $f(X)_t$ depends only on
the input at $t$ (and fixed metadata such as $t,T$), not on other positions.
\end{proposition}
\begin{proof}
By left causality, $f(X)_t$ depends only on the prefix through $t$.
Equivariance expresses this value as an invertible channel transform of
$f(\mathcal RX)_{T+1-t}$. Left causality on the reversed input makes that
quantity depend only on the original suffix from $t$ onward. Now take two
inputs that agree at $t$. Form a third using the first input's prefix and
the second input's suffix; it exists because the domain is a full product.
Prefix dependence equates its output with the first input, and suffix
dependence equates it with the second. Hence only the shared token can matter.
\end{proof}

The assumptions matter. This is a theorem about same-position
sequence-to-sequence outputs at fixed length, not a prohibition on all
strand-aware causal models. The paired representation avoids the conclusion
by being bidirectional, not left-causal. A forward-only branch can be used for
causal decoding, but no longer provides the full paired-output guarantee.
Record-level invariant outputs emitted after the record ends are another
different contract.

For autoregressive training, a forward conditional
$p(x_{t+1}\mid x_{\leq t})$ and a reverse-oriented conditional are different
prediction problems. One may train both directions with shared weights and
correctly transformed targets; this does not justify averaging their aligned
logits while evaluating a forward conditional. A target mask must also be
transformed with coordinates. Bidirectional masked-token objectives require
masking the target information before either pass sees it.

\section{Reverse the record, not the transport chunks}

Let a complete record be partitioned as $X=C_1C_2\cdots C_m$.
Reverse complement is an anti-homomorphism for concatenation:
\begin{equation}
 \mathcal R(C_1C_2\cdots C_m)
 =\mathcal R(C_m)\cdots\mathcal R(C_2)\mathcal R(C_1).
 \label{eq:chunks}
\end{equation}
The order of chunks and the coordinates within each chunk both reverse.
For ACG$\mid$TA$\mid$CC, the correct opposite record is
GG$\mid$TA$\mid$CGT. Independently transforming chunks in arrival order gives
CGT$\mid$TA$\mid$GG, a different record.

\Plate{05-chunks}{Chunk boundaries are a transport partition, not a new
biological sequence order. The reverse pass starts with the final chunk and
continues its state into the preceding chunk. Each record has its own pair of
boundary states; concatenating unrelated records does not license state reuse.}{fig:chunks}

\Listing{chunked_dual}

The forward pass carries its final state left to right. The opposite pass
starts from zero at the opposite record boundary, visits chunks in reverse
order and then realigns the concatenated outputs. Exact-arithmetic parity with
Equation~\ref{eq:lift} follows because the recurrence sees the same token order
and the same state at every step. Tests check every three-way partition of
several record lengths, including empty chunks, and compare parameter and
input gradients as well as forward values.

This implementation is offline and retains output histories. In a real
system, reverse traversal requires an available complete record, a reverse
data source, or stored input that can be reread. It is not supplied by a
constant-size forward state in general. For a scalar recurrence,
the feature sequences $(1,-a)$ and $(0,0)$ have the same forward final state
zero, but their reverse final states are $1-a^2$ and zero. This counterexample
concerns the real-valued feature domain; it is not an assertion that these
numbers are literal nucleotide tokens.

Padding requires similar care. Our reference processes unpadded records.
If right padding is reversed along with the record, it becomes left padding.
Advancing the recurrence on a zero row decays the state and is not the same
as skipping a padding slot. A production implementation should transform the
validity mask and obey a declared hold/advance policy. Likewise, a cached
state needs record identity, orientation, model revision and boundary location,
not just the correct vector dimension.

\section{Differentiate the shared computation}

Sharing weights means gradients add across both uses. If
$E_\theta(X)=[F_\theta(X)\mid R_TF_\theta(\mathcal RX)]$ and a loss has branch
cotangents $(g_f,g_b)$, then schematically
\begin{equation}
 \nabla_\theta\mathcal L
 =(\partial_\theta F_\theta(X))^\top g_f+
  (\partial_\theta F_\theta(\mathcal RX))^\top R_Tg_b .
\end{equation}
There is one parameter object, not two weights periodically copied after
optimization. Position reversal and complement are permutations, whose
adjoints are their inverses; in this case they are involutions.
Consequently, gradients through the transformed input return to the
corresponding original coordinates.

For a loss invariant under the output action and consistently transformed
targets, the losses on $X$ and $\mathcal RX$ are equal in exact arithmetic.
Their parameter gradients are equal wherever differentiable. This does not
mean the two branch contributions within one evaluation are equal.
Finite differences, automatic-differentiation parity, and deliberately broken
weight sharing form separate checks. Our tests include all three forms of
reasoning: a closed-form recurrence oracle, numerical gradient checking, and
negative examples whose symmetry residual is nonzero.

\section{A falsifiable strand hypothesis, not a biological shortcut}

Caduceus provides a primary example of distinguishing bidirectionality,
shared reverse-complement processing, equivariant embeddings and output
heads~\cite{caduceus}. It also separates architectural equivariance from
post-hoc conjoining. Our paired-block action and simple recurrence are an
independently specified teaching reference, not a reproduction of its
MambaDNA implementation or reported benchmarks.

For DOGMA, formulate the scientific claim narrowly: does a declared strand
constraint improve a specific task under matched budgets and correct labels?
Compare a forward baseline, an unconstrained bidirectional baseline,
augmentation, the shared dual lift and an invariant readout only when the
task warrants invariance. Match or report parameters, FLOPs, input access,
state storage and preprocessing. Sharing parameters does not eliminate the
arithmetic of a second traversal.

Split data before orientation augmentation. Otherwise a record can appear in
training while its reverse complement appears in testing. Group related loci
and homologous examples according to the scientific question; exact RC
deduplication alone does not remove biological relatedness. A palindromic
record satisfies $\mathcal RX=X$, so its paired output must satisfy
$E(X)=\mathcal JE(X)$. That relates opposite positions and branches, not
necessarily the two branches at the same position.

Report task metrics alongside symmetry residuals and failure cases. Zero
symmetry residual proves only the implementation's transformation contract.
A constant predictor can satisfy it perfectly and still be useless. Conversely,
a task with deliberately oriented labels can be harmed by an invariant head
that erases needed distinctions. The output action is a scientific choice
to test, not an architectural slogan.

\section{Exercises with worked answers}
\begin{enumerate}
\item Compute the reverse complement of ACGA. \textbf{Answer:} Complement
TGCT, then reverse to TCGT; complement-only is not the required reading order.
\item Transform $[2,5)$ in a record of length $9$. \textbf{Answer:}
$[9-5,9-2)=[4,7)$. Applying the rule again returns $[2,5)$.
\item Prove $\mathcal R^2=I$. \textbf{Answer:} Row reversal and channel
complement commute, and each applied twice is identity.
\item What output action fits two strand-specific tracks? \textbf{Answer:}
Reverse positions and swap track labels, provided the assay metadata
transforms in the same way.
\item Why is the second branch reversed after evaluation? \textbf{Answer:}
Its row $T+1-t$ corresponds to original position $t$; without realignment
the concatenated rows describe different positions.
\item Derive $h_3$ for projected inputs $(1,2,3)$ and $a=1/2$.
\textbf{Answer:} $h_1=1$, $h_2=2.5$, $h_3=4.25$.
\item Solve the block constraints for a linear paired readout.
\textbf{Answer:} $MJ=J'M$ implies $A=D_0$ and $B=C$; an affine
bias must have equal paired blocks.
\item Why does equivariance not permit causal use of the entire dual state?
\textbf{Answer:} Its aligned reverse branch uses the suffix. Symmetry
does not remove those dependencies.
\item Reverse complement ACG$\mid$TA$\mid$CC. \textbf{Answer:}
GG$\mid$TA$\mid$CGT; changing each chunk without reversing their order is wrong.
\item Does a zero padding row act as identity? \textbf{Answer:}
Not in Equation~\ref{eq:recurrence}: it maps $h$ to $a\odot h$.
Skipping an invalid slot needs an explicit mask rule.
\item Must a palindromic input have equal branch features at each position?
\textbf{Answer:} No. It requires $u_t=v_{T+1-t}$; equality at $t$
only follows where additional conditions, such as a central fixed position, apply.
\item Design a fair strand-ablation study. \textbf{Rubric:} Specify label
transformations and orientation metadata; split related records before
augmentation; compare forward, bidirectional and constrained variants;
match or report computational budgets; test symmetry and task accuracy
separately; include an oriented task and a constant-predictor control.
\end{enumerate}

\clearpage
\section{Vocabulary and the integration boundary}
\Term{Reverse complement}: reverse sequence order and complement base
identities. \Term{Output action}: the specified transformation of predictions
when the input representation changes. \Term{Parameter tying}: multiple
computational uses of one parameter object. \Term{Bidirectionality}: use of
both prefix and suffix context. \Term{Anti-homomorphism}: a transformation
that reverses the order of composition, as in Equation~\ref{eq:chunks}.

Chapter 22 will assemble the DOGMA reference research model. Its first job is
not to add more biological names: it is to reconcile the token, regulation,
memory, strand, objective and execution contracts into a model whose successes
and failures can be measured without changing the claim halfway through.
